% !TeX root = ../main.tex
%
% Licensing and attribution statement.
%
% Every licence named here was verified from the asset itself (the package
% metadata in the pinned environment, or the repository's LICENSE file), not
% recalled.  Where no licence declaration could be established, this file says
% so rather than asserting one -- the checklist guidelines explicitly permit
% "we were unable to find the license".
%
% ANONYMITY: the code licence is named without its copyright line, because the
% copyright line contains the author's name.  Do not paste the LICENSE header.

\section{Licensing and attribution of existing assets}
\label{sec:stmt:licensing}

\paragraph{Dataset.}
All experiments use CIFAR-10 \citep{krizhevsky2009cifar}, obtained through the
standard \texttt{torchvision} dataset loader and used unmodified. We credit it
by citing the technical report that released it, which is the attribution its
authors ask for. We were not able to establish a formal licence declaration
accompanying the CIFAR-10 distribution, and we state that rather than assert a
licence we cannot evidence; the dataset is distributed publicly for research
use, and our use --- non-commercial academic benchmarking --- is within that
purpose.

\paragraph{Pre-trained weights.}
Every dense model in this paper is initialised from ImageNet-pretrained
\texttt{torchvision} classification weights rather than from random
initialisation, and the pruned arms are initialised from the resulting dense
checkpoint. Those weights were trained on ImageNet
\citep{deng2009imagenet}, whose own terms of access grant use for
non-commercial research and educational purposes; we use them within that
scope and release no derived weights. The weights are distributed as part of
\texttt{torchvision} and carry its licence, below.

\paragraph{Architectures.}
The backbones are published architectures, not ours: the residual networks of
\citet{he2016deep} and the VGG family of \citet{simonyan2015very}, taken from
their \texttt{torchvision} implementations. The training objective this paper
studies is the module decomposition of CAP \citep{xu2022cap}; the paper states
plainly that the decomposition is theirs and identifies precisely where our
implementation departs from their published form. The sparsity schedule is that
of \citet{zhu2018prune} and the prune-and-retrain paradigm that of
\citet{han2015learning}. The pruning criteria are those of \citet{lee2019snip}
and \citet{verdenius2020snipit} for the sensitivity score and its iterative
re-application, and \citet{sun2024wanda} for the activation-weighted score.

\paragraph{Software.}
The environment is pinned. \texttt{PyTorch} 2.5.1 and \texttt{torchvision}
0.20.1 are used under the BSD 3-Clause licence. The codebase also depends on
\texttt{transformers} 4.46.3 and \texttt{datasets} 3.1.0, both under the
Apache License 2.0; these support the language-model configurations that the
framework implements but that this paper does not use --- no result reported
here comes from a HuggingFace checkpoint, and if a text result is added, the
licence on each model card must be checked individually, since the checkpoints
in that family do not share a single licence.

\paragraph{Our own code.}
The training and analysis code written for this project is licensed under the
MIT License. It is not linked from this submission, for the anonymity reason
given in the reproducibility statement, and the licence text is stated here
without its copyright line for the same reason.

\paragraph{Published numbers quoted for comparison.}
The comparator accuracies this paper places beside its own are quoted from the
tables of the papers that reported them, with the source and the training recipe
named in each case, and are not re-run by us. They are used as published
reference points under normal scholarly citation, and the paper states the
difference in training budget that makes a direct comparison of absolute
accuracies inappropriate.
